% system-model.tex -- Sequencer Paper (BNR-2026-009)

\section{System Model}

\subsection{Target Architecture}

The Basis Network enterprise zkEVM L2 operates a per-enterprise chain architecture
(design decision TD-005 from the foundational specifications). Each enterprise deploys
an independent L2 instance consisting of:

\begin{itemize}
  \item A \textbf{Go-based sequencer node} that receives transactions from enterprise
    applications (PLASMA, Trace), orders them via FIFO policy, produces L2 blocks at
    1-second intervals, and submits batches to L1.
  \item A \textbf{Rust-based ZK prover} that generates validity proofs for each batch
    of state transitions.
  \item \textbf{L1 smart contracts} on Basis Network (Avalanche Subnet-EVM, Chain~ID
    43199) that verify proofs and anchor state roots.
\end{itemize}

The sequencer is the focus of this paper. It sits between the transaction ingestion
layer (enterprise applications submitting via RPC) and the proving layer (Rust prover
consuming state transitions).

\subsection{Design Decisions}

Table~\ref{tab:design-decisions} summarizes the architectural choices for the Basis
Network sequencer and their rationale.

\begin{table}[htbp]
\caption{Sequencer Design Decisions}
\label{tab:design-decisions}
\centering
\begin{tabular}{lll}
\toprule
\textbf{Decision} & \textbf{Choice} & \textbf{Rationale} \\
\midrule
Sequencer type & Single-operator & Per-enterprise chains \\
Ordering policy & FIFO (arrival) & Zero-fee, no MEV \\
Block time & 1~second & Enterprise latency req. \\
Block gas limit & 10M (configurable) & Bounds execution/block \\
Forced inclusion & Arbitrum-style & FIFO queue, 24h deadline \\
Batch strategy & Hybrid (size OR time) & Validated in RU-V4 \\
Mempool & Simple FIFO queue & No priority needed \\
\bottomrule
\end{tabular}
\end{table}

\subsection{Sequencer State Machine}

The sequencer operates as a deterministic state machine with the following cycle,
repeated every block interval $\Delta t = 1$~second:

\begin{enumerate}
  \item \textbf{Drain forced inclusion queue}: Check for L1-submitted transactions
    whose deadline has expired. These are prepended to the block with highest priority.
  \item \textbf{Drain mempool}: Extract up to $N_{max}$ transactions from the FIFO
    mempool, preserving arrival order.
  \item \textbf{Assemble block}: Combine forced transactions and mempool transactions
    into an L2 block with sequential numbering, parent hash linkage, and timestamp.
  \item \textbf{Advance state}: Update block height, parent hash, and metrics counters.
\end{enumerate}

The forced inclusion queue takes absolute priority: forced transactions are always
included before mempool transactions in each block. This guarantees that censored
transactions bypass the sequencer's mempool entirely.

\subsection{Formal Properties}

The sequencer must satisfy the following properties:

\textbf{FIFO Ordering}: For any two transactions $tx_a$ and $tx_b$ where $tx_a$
arrives before $tx_b$ and both are included in block $B_k$, $tx_a$ must appear
before $tx_b$ in $B_k$'s transaction list:
\[
  \text{arrival}(tx_a) < \text{arrival}(tx_b) \implies \text{index}(tx_a, B_k) < \text{index}(tx_b, B_k)
\]

\textbf{Forced Inclusion Liveness}: For any transaction $tx_f$ submitted to the
forced inclusion queue at time $t_0$ with deadline $d$, there exists a block $B_k$
with timestamp $t_k \leq t_0 + d$ such that $tx_f \in B_k$:
\[
  \forall tx_f \in Q_{forced}: \exists B_k \text{ s.t. } tx_f \in B_k \land t_k \leq t_0 + d
\]

\textbf{Block Production Regularity}: The sequencer produces a block at every tick
interval $\Delta t$, even if no transactions are pending (empty blocks are permitted
but configurable):
\[
  \forall k \geq 1: t_{B_k} - t_{B_{k-1}} \leq \Delta t + \epsilon
\]
where $\epsilon$ accounts for processing overhead.

\textbf{Mempool Capacity Enforcement}: The mempool has a maximum capacity $C_{max}$.
When full, new insertions are rejected (not silently dropped):
\[
  |Q_{mempool}| = C_{max} \implies \text{Insert}(tx) = \text{false}
\]
